\section{Analysis}
\label{sec:analysis}

We probe the trained hormone-routing variants to characterise what the router learns despite its lack of measurable contribution to the loss. Three findings drive the discussion: (i) every variant develops semantically-coherent per-genre specialisation; (ii) the upper attention layer engages the gate while the lower layer does not; (iii) the asymmetric engagement holds across all unforced variants and is structurally similar to what \textsc{HR-noext} produces despite its loss damage.

\subsection{The router learns genre-conditional specialisation}
\label{subsec:routing-specialisation}

We probe the trained \textsc{HR}, \textsc{HR-rand}, \textsc{HR-fixedgate}, and \textsc{HR-noext} models by recording, for each of six text genres (news, fiction, code, dialogue, technical exposition, urgent / emergency text), the average of the router's softmax output over all token positions in a representative passage. Results for the upper attention layer (the engaged one in all four variants) are shown in Figure~\ref{fig:routing}; the analogous distributions for the lower attention layer follow the same pattern in muted form. We summarise specialisation strength via the maximum pairwise $L_1$ distance between the six genre distributions: $0$ corresponds to identical routing across genres (no specialisation), $2$ to fully disjoint distributions (perfect one-hot specialisation).

\begin{figure}[h]
\centering
\TBD{Insert fig4\_routing\_heatmaps\_layer11.pdf — three side-by-side heatmaps for HR, HR-rand, HR-fixedgate; rows = genres, columns = hormone indices, cells annotated with router weights.}
\caption{Per-genre router distributions at the upper attention layer (layer $11$), one heatmap per variant. The router learns a strongly genre-discriminative distribution under all four regimes. \textsc{HR} (max pairwise $L_1 = 0.60$) routes news and tech text preferentially to the urgency direction (ADR), fiction to curiosity (CDO), code to skepticism (LCO), urgent text to warmth (OXY) --- consistent with the semantic labels of the extracted directions. \textsc{HR-rand} ($L_1 = 1.37$), \textsc{HR-fixedgate} ($L_1 = 1.56$), and \textsc{HR-noext} (avg-layer $L_1 = 0.73$, layer-$11$ alone $\approx 1.4$) develop equally or more specialised distributions, but the assignments are arbitrary --- the directions in those banks have no semantic meaning. The specialisation is therefore a property of the routing architecture and the data distribution, not of the directions. Crucially, the variant with the \emph{highest} specialisation strength (\textsc{HR-fixedgate}) achieves the same val PPL as the variant with the \emph{lowest} (\textsc{HR}) and as the unmodified Base. Interpretive specialisation strength does not predict downstream loss.}
\label{fig:routing}
\end{figure}

\paragraph{Interpretive note: \textsc{HR}'s specialisation is interpretable; the others are not.} In \textsc{HR} the directions were extracted via contrast pairs labelled by intended affective dimension (e.g.\ urgency, curiosity, skepticism). The router's preferential routing of news $\rightarrow$ ADR, fiction $\rightarrow$ CDO, code $\rightarrow$ LCO is consistent with intuitive readings of those genres. In \textsc{HR-rand} and \textsc{HR-fixedgate} the same pattern of strong genre-specific preference emerges, but the chosen directions are arbitrary samples from the residual subspace --- their indices carry no meaning beyond ``the third unit vector,'' ``the fifth unit vector.'' The structural similarity of the specialisations across these very different banks is the load-bearing observation: the architecture+data system produces routing organisation regardless of the bank's content. Extracted directions provide interpretive labels for the resulting clusters; they do not produce the clusters.

\paragraph{\textsc{HR-noext} also specialises.} The variant with the worst validation perplexity also develops strong specialisation: at the upper attention layer max pairwise $L_1 = 0.73$ (averaged across both layers; layer-$11$ alone is higher), with heavy preference for the LCO and SELF indices. Its loss damage cannot therefore be attributed to a failure to specialise.

\subsection{Layer-asymmetric engagement}

The output gate $g$ at the upper attention layer (index $11$) opens during Stage-$2$ training in every variant: $\textsc{HR}: g = +0.231$; $\textsc{HR-rand}: g = -0.122$; $\textsc{HR-fixedgate}: g = 1.0$ (frozen); $\textsc{HR-noext}: g = +0.366$. The lower attention layer (index $5$) stays effectively closed in the unforced variants ($\textsc{HR}: g = -0.0002$; $\textsc{HR-rand}: g = +0.0007$; $\textsc{HR-noext}: g = +0.012$). This is consistent with the residual-stream organisation usually observed in transformers~\citep{zou2023representation}: the upper layers carry more decision-relevant signal and are therefore where the router's perturbation has somewhere to go.

\paragraph{Magnitude trajectories.} The per-hormone learned magnitudes $m_i$ in \textsc{HR} were initialised at $1.0$. By the end of training the upper-layer magnitudes had grown modestly (range $1.06$--$1.14$) while the lower-layer magnitudes remained essentially at initialisation ($0.97$--$0.98$), consistent with the asymmetric gate engagement: lower-layer magnitudes received no gradient signal because their contribution was masked by the closed gate.

\subsection{Discussion: organisation independent of objective}

The empirical pattern --- specialisation that develops across all bank choices, gate engagement that is non-trivial but small in magnitude, validation loss invariant to all of this --- supports a non-trivial claim that we now state directly.

\paragraph{Specialisation strength does not correlate with loss.} Across the four hormone-routing variants, the variant with the \emph{highest} specialisation (\textsc{HR-fixedgate}, $L_1 = 1.56$) is one of the four that converge to identical val PPL with Base; the variant with the \emph{lowest} specialisation among the well-behaved ones (\textsc{HR}, $L_1 = 0.60$) likewise matches Base; and the variant with intermediate specialisation that does change loss (\textsc{HR-noext}, $L_1 \approx 0.73$) changes it in the wrong direction. There is no monotone relationship between how much the router specialises and whether the model's predictions improve. Interpretive metrics --- the kind of metrics typically used to claim a mechanism is doing something useful --- are decoupled from the objective the mechanism is trained against. Two complementary readings of this decoupling are admissible:

\begin{description}
    \item[Structural.] The routing module is a parameter-efficient channel that the architecture+data system uses to organise its representations regardless of whether that organisation is loss-relevant. The gradient that reaches the router weights is sufficient to produce a discriminative distribution; the gradient that reaches the gate, after the perturbation has been absorbed by downstream layers, is not sufficient to push the gate into a regime where the perturbation reshapes predictions.
    \item[Capacity-bottlenecked.] At $125$M parameters and $\dmodel = 1024$ the residual stream is too low-dimensional for a small additive perturbation along seven fixed directions to reliably encode loss-reducing information; the same mechanism at larger $\dmodel$ might find the gate value where the perturbation begins to matter.
\end{description}

The forced-engagement variant (Section~\ref{subsec:hrforced}) provides partial discrimination: \TBD{narrative depending on whether forced engagement changes loss}.

\paragraph{Content-invariance versus timing-sensitivity.} A separate way of reading the same data is the dimensional axis along which the model is invariant. The model is invariant to the routing module's \emph{content} --- swapping the extracted bank for random unit vectors (\textsc{HR-rand}) or freezing the gate at $1.0$ (\textsc{HR-fixedgate}) does not perturb val PPL. The model is \emph{not} invariant to the module's \emph{training-schedule timing} --- introducing the routing during base pretraining (\textsc{HR-noext}) does perturb val PPL, and harmfully. The architectural homeostasis we observe is a property of a model that already has a stabilised representation; it is not a property of the architecture-plus-data system from a fresh initialisation. This pattern is consistent with a picture in which the model develops, over Stage~1, the representational structure required to absorb additive perturbations as background noise; the same perturbation introduced before this structure is in place becomes a moving-target distractor that interferes with the formation of the structure itself.

\paragraph{Bounded mechanism as a design property.} A weaker but tangible claim survives regardless of which interpretation of the loss-decoupling is correct: the mechanism is \emph{bounded}. With proper initialisation, hormone routing cannot harm performance --- random vectors and forced-magnitude gates both yield base-level loss. This is a useful design property in its own right. Most architectural additions to language models are introduced with a presumed positive contribution and a worry about unintended side-effects; hormone routing inverts the picture, with no positive contribution observable at this scale but no side-effects either. Combined with the routing-specialisation finding, this means the mechanism is well-suited to deployment as an interpretability or steering substrate that can be inspected during training without affecting the loss surface.

\subsection{Limitations}

\paragraph{Scale.} Our experiments are at the $125$M-parameter scale on a $2$B-token budget. Whether the routing-organisation-without-loss-relevance finding persists, intensifies, or reverses at larger scale is unknown. Scaling laws~\citep{hoffmann2022training} predict that the relative impact of architectural modifications can change with model size, and in particular the ``capacity-bottlenecked'' interpretation of our null val-PPL result becomes testable only at scales where $\dmodel$ is large enough to provide meaningful expressive room for the perturbation.

\paragraph{Pretraining-only.} We do not perform supervised fine-tuning or RLHF. The hormone routing mechanism is in principle compatible with these stages --- the routing weights and magnitudes can be further updated, and re-extraction can target newly-relevant behavioural axes --- but we leave that investigation to future work.

\paragraph{Single-architecture study.} Both contributions are designed for the Mamba/MQA hybrid family. Intracellular attention generalises naturally to any SSM with a state shape that admits the slot interpretation (S4, S5, Mamba-2). Hormone routing is architecture-agnostic and could be applied to a pure-transformer base, but we have not validated this empirically.

\paragraph{Reference-scan throughput for IA.} Full-budget training of the IA variants requires a fused CUDA kernel to remove the per-step Python loop overhead of the reference scan. The architectural and parameter-cost analysis we provide is independent of the kernel, but the empirical case for IA at the full budget is not.
